% Source PDF pages 110--113; manual pages 2-81--2-84.
\subsection{Control Variable Solution Process}\label{sec:control-variable-solution}

TAOS separates guidance rules that directly specify control variables from those
that do not. It compares the directly specified controls with
\cref{tab:body-attitude-control-variables} and selects the lowest-numbered set
that contains them. If no set is consistent with the input, TAOS reports an input
error. Direct controls are assigned their input values; the remaining controls and
indirect guidance rules are retained for solution.

The indirect solution is iterative. The procedure, shown in
\cref{fig:guidance-loop}, is the \emph{guidance loop} and forms part of each
equations-of-motion derivative evaluation.

\begin{figure}[htbp]
  \centering
  \begin{tikzpicture}[
  >=Latex,
  node distance=0.55cm and 1.05cm,
  box/.style={draw,minimum width=3.15cm,minimum height=0.75cm,align=center},
  sidebox/.style={draw,text width=3.25cm,minimum height=2.9cm,align=center},
  every path/.style={line width=0.7pt}
]
  \node[box] (forces) {Calculate Forces};
  \node[box,below=of forces] (accel) {Calculate Accelerations};
  \node[box,below=of accel] (rates) {Calculate Rates};
  \node[box,below=of rates] (loop) {Guidance Loop};
  \node[sidebox,right=1.55cm of accel] (vary)
    {Vary control variables to get accelerations required by guidance rules};

  \draw[->] ($(forces.north)+(0,0.55)$) -- (forces.north);
  \draw[->] (forces) -- (accel);
  \draw[->] (accel) -- (rates);
  \draw[<->] (rates) -- (loop);
  \draw[->] (loop.east) -| (vary.south);
  \draw[->] (vary.north) |- ($(forces.north)+(0,0.55)$);
  \draw[->] (loop.south) -- ++(0,-0.65)
    node[below,align=center] {Solution\\complete};
\end{tikzpicture}

  \caption{Guidance loop.}
  \label{fig:guidance-loop}
\end{figure}

Because acceleration may change instantaneously along the trajectory, each indirect
guidance rule is associated with the solution variable listed in
\cref{tab:indirect-guidance-rules}. The rule supplies a desired value, and TAOS
varies a corresponding control variable until the computed and desired solution
variables agree. The resulting nonlinear system is solved by a multidimensional
Newton--Raphson method.

For quantities such as $C_L$ or an axial specific load, the desired solution variable
is given directly. For altitude and flight-path-angle rules, the input specifies an
integral or double integral of the solution variable. TAOS converts that command to
a desired instantaneous acceleration by fitting a second- or third-order transition
in time.

\taossubhead{Parabolic Transition}

Consider the command to fly at a constant geodetic vertical flight-path angle of
$5^\circ$. It specifies a desired state $\gamma_{gd}'=5^\circ$ and desired rate
$\dot\gamma_{gd}'=0$. The prime denotes a desired value. TAOS assumes that
$\gamma_{gd}$ follows a second-order parabolic transition from its current value to
the desired state over the user-supplied interval $\Delta t_{\mathrm{guid}}$, as
shown in \cref{fig:parabolic-guidance-transition}.

\begin{figure}[htbp]
  \centering
  \begin{tikzpicture}[>=Latex,x=1.05cm,y=1.05cm,line cap=round]
  \draw[->] (0,0) -- (8.7,0);
  \draw[->] (0,-0.25) -- (0,3.3);
  \draw[line width=0.9pt]
    plot[smooth] coordinates {(0,2.75) (0.8,2.0) (2.0,1.28) (3.6,0.82) (5.5,0.42) (7.7,0.04)};
  \draw (0,0.05) -- (8.1,0.05);
  \fill (0,2.75) circle (2pt);
  \fill (7.7,0.05) circle (2pt);
  \draw[<->] (0,-0.45) -- node[fill=white,inner sep=1pt] {$\Delta t_{\mathrm{guid}}$} (7.7,-0.45);
  \draw (0,-0.18) -- (0,-0.62);
  \draw (7.7,-0.18) -- (7.7,-0.62);
  \node[left] at (0,2.75) {$\gamma_{gd}$};
  \node[left] at (0,0.05) {$\gamma_{gd}'$};
  \node[below] at (0,-0.62) {$t$};
  \node[below] at (7.7,-0.62) {$t+\Delta t_{\mathrm{guid}}$};
  \node[above right,align=center] at (0.15,2.72) {Current\\state};
  \node[above left,align=center] at (7.62,0.2) {Desired\\state};
  \draw[->] (4.85,2.15) -- (4.15,0.67);
  \node[above,align=center] at (4.85,2.15) {Parabolic function\\of time};
\end{tikzpicture}

  \caption{Parabolic transition to a desired state.}
  \label{fig:parabolic-guidance-transition}
\end{figure}

The interval $\Delta t_{\mathrm{guid}}$ acts as a time constant. Large values produce
small corrections and slow convergence to the desired condition. Small values
produce rapid corrections, but values that are too small can cause oscillation or
divergence.

The desired acceleration $\dot\gamma_{gd}'$ at the current time is obtained from a
parabolic curve fitted through the current and desired conditions:
\begin{equation}
  \gamma_{gd}'(t)=a t^2+b t+c,
  \label{eq:parabolic-guidance-state}
\end{equation}
whose derivative is
\begin{equation}
  \dot\gamma_{gd}'(t)=2at+b.
  \label{eq:parabolic-guidance-rate}
\end{equation}
The coefficients used by TAOS are
\begin{equation}
  a=\frac{\gamma_{gd}-\gamma_{gd}'
    +\dot\gamma_{gd}'\,\Delta t_{\mathrm{guid}}}
    {\Delta t_{\mathrm{guid}}^2},
  \label{eq:parabolic-guidance-a}
\end{equation}
\begin{equation}
  b=\dot\gamma_{gd}'-2a\left(t+\Delta t_{\mathrm{guid}}\right).
  \label{eq:parabolic-guidance-b}
\end{equation}
The function \nolinkurl{parab_guidance_correction} in the \texttt{derivs} module
applies this procedure to single-derivative solution variables such as
$\dot\gamma_{gd}$, $\dot\psi_{gd}$, and $\lVert\dot{\vec V}\rVert$.

\taossubhead{Cubic Transition}

For double-derivative solution variables such as $\ddot h$, TAOS fits a third-order
curve with \nolinkurl{cubic_guidance_correction}, as illustrated in
\cref{fig:cubic-guidance-transition}. Current and desired altitude and altitude rate
provide four boundary conditions.

\begin{figure}[htbp]
  \centering
  \begin{tikzpicture}[>=Latex,x=1.0cm,y=1.0cm,line cap=round]
  \draw[->] (0,0) -- (8.7,0);
  \draw[->] (0,-0.25) -- (0,3.25);
  \draw[line width=0.9pt]
    plot[smooth] coordinates {(0,2.25) (0.8,2.75) (1.8,2.45) (3.0,1.45) (4.5,0.68) (6.2,0.25) (7.7,0.05)};
  \draw (0,0.05) -- (8.1,0.05);
  \fill (0,2.25) circle (2pt);
  \fill (7.7,0.05) circle (2pt);
  \draw[<->] (0,-0.45) -- node[fill=white,inner sep=1pt] {$\Delta t_{\mathrm{guid}}$} (7.7,-0.45);
  \draw (0,-0.18) -- (0,-0.62);
  \draw (7.7,-0.18) -- (7.7,-0.62);
  \node[left,align=right] at (0,2.25) {$h$ and $\dot h$};
  \node[left,align=right] at (0,0.05) {$h'$ and $\dot h'$};
  \node[below] at (0,-0.62) {$t_1=t$};
  \node[below] at (7.7,-0.62) {$t_2=t+\Delta t_{\mathrm{guid}}$};
  \node[above right,align=center] at (0.15,2.35) {Current\\state};
  \node[above left,align=center] at (7.6,0.18) {Desired\\state};
  \draw[->] (4.65,2.25) -- (4.0,0.92);
  \node[above,align=center] at (4.65,2.25) {Cubic function\\of time};
\end{tikzpicture}

  \caption{Cubic transition to a desired state.}
  \label{fig:cubic-guidance-transition}
\end{figure}

The transition is represented by
\begin{equation}
  h'(t)=a t^3+b t^2+c t+d,
  \label{eq:cubic-guidance-state}
\end{equation}
and differentiation twice gives
\begin{equation}
  \ddot h'(t)=6at+2b.
  \label{eq:cubic-guidance-acceleration}
\end{equation}
With $t_1=t$ and $t_2=t+\Delta t_{\mathrm{guid}}$, the printed TAOS formulation
uses
\begin{equation}
\begin{aligned}
 b={}&\Bigl[(h-h')(t_2^2-t_1^2)
   +\dot h(t_2-t_1)(t_2^2-t_1^2)\\
 &\quad+(\dot h'-\dot h)
   \left(\frac{2}{3}t_1^3-t_1^2t_2+\frac{1}{3}t_2^3\right)\Bigr]\\
 &\Bigg/\Bigl[2(t_2-t_1)
   \left(\frac{2}{3}t_1^3-t_1^2t_2+\frac{1}{3}t_2^3\right)
   -(t_2^2-t_1^2)(t_1^2-2t_1t_2+t_2^2)\Bigr],
\end{aligned}
\label{eq:cubic-guidance-b}
\end{equation}
\begin{equation}
  a=\frac{\dot h'-\dot h-2b(t_2-t_1)}{3(t_2^2-t_1^2)}.
  \label{eq:cubic-guidance-a}
\end{equation}

\taossubhead{Newton--Raphson Solution}

Let the free control variables be collected in $\vec x$, and let the residual vector
$\vec f$ contain the differences between desired and computed solution variables,
$f_i=f_i'-f_i$. The multidimensional Newton--Raphson iteration is
\begin{equation}
  \vec x_{n+1}=\vec x_n-[J]^{-1}\vec f,
  \label{eq:guidance-newton-update}
\end{equation}
where the source defines the Jacobian as
\begin{equation}
  [J]=
  \begin{bmatrix}
    \dfrac{\partial f_1}{\partial x_1} &
    \dfrac{\partial f_2}{\partial x_1} &
    \dfrac{\partial f_3}{\partial x_1} & \cdots\\[0.9em]
    \dfrac{\partial f_1}{\partial x_2} &
    \dfrac{\partial f_2}{\partial x_2} &
    \dfrac{\partial f_3}{\partial x_2} & \cdots\\[0.9em]
    \dfrac{\partial f_1}{\partial x_3} &
    \dfrac{\partial f_2}{\partial x_3} &
    \dfrac{\partial f_3}{\partial x_3} & \cdots\\
    \vdots & \vdots & \vdots & \ddots
  \end{bmatrix}.
  \label{eq:guidance-jacobian}
\end{equation}
The Jacobian and residual are evaluated at $\vec x_n$. Defining
$\vec z_{n+1}=\vec x_{n+1}-\vec x_n$ yields the linear system
\begin{equation}
  [J]\vec z_{n+1}=-\vec f,
  \label{eq:guidance-gaussian-system}
\end{equation}
which TAOS solves by Gaussian elimination rather than explicit matrix inversion.

For the first point in a trajectory, each free-control initial estimate is placed midway
between its upper and lower limits. At subsequent points, the preceding solution is
used as the next initial estimate, relying on the usual assumption that controls vary
smoothly along the trajectory.
